I study the effect of electricity prices on market entry of electric vehicle (EV) charging stations. I compile a novel dataset linking commercial electricity prices to charging stations across U.S. ZIP codes from 2015–2022. I convert complex pricing schedules into a standardized cost metric by constructing station-level load profiles from observed charging sessions and calculating monthly bills. Using synthetic control and local projection difference-in-differences designs, I estimate the effects of new price schedules for EV charging stations, aimed at reducing high demand charges.
The results show that these schedules substantially lowered demand charges for DC fast charging stations by about 50\% and increased the entry of DC fast charging ports by about 50\% per zip code, underscoring the role of targeted rate design in accelerating EV infrastructure growth.


%This result is consistent with a reduction in total electricity charges, driven by significant decreases in energy and demand charges, 54\% and 50\%, respectively, underscoring the role of targeted rate design in accelerating EV infrastructure growth.

% Shorter version
% This paper examines how electricity rate plans affect EV charging station installations, using a novel dataset I construct that combines commercial rates and charging stations in U.S. zip codes from 2015 to 2022. I aggregate intricate pricing into a single cost metric via bill simulation. Using synthetic control and local projection difference-in-differences, I estimate how new electricity tariffs aimed at charging station operators—designed to alleviate high demand charges—affected their electricity bills and station entries. These specialized rates led to lower demand charges and 1-2 additional charging ports per zip code, highlighting the effectiveness of targeted rate design in promoting EV infrastructure development.